% 数学ノート集・共通テンプレート v2（積分論2026年度第1部の版面を参考）
% A5判。章内で定義・定理・例・注意等の番号を共通化。
% コンパイル： platex -interaction=nonstopmode math_note_template_v2.tex (2回)
%              dvipdfmx math_note_template_v2.dvi
\documentclass[a5paper,10pt,onecolumn]{jsarticle}
\usepackage[dvipdfmx,top=22mm,bottom=20mm,inner=19mm,outer=18mm,headheight=13pt,headsep=6mm]{geometry}
\usepackage{amsmath,amssymb,amsthm,mathtools}
\usepackage{enumitem,fancyhdr,titlesec}
\usepackage[dvipdfmx,bookmarks=true,colorlinks=false]{hyperref}
\usepackage{pxjahyper}

% ----- 差し替える文書情報 -----
\newcommand{\NoteTitle}{自由対象と普遍性}
\newcommand{\NoteSeries}{数学ノート集}

% ----- 文字・組版 -----
% 参考PDFの質素な明朝＋Computer Modernの印象を引き継ぎながら，
% A5判で窮屈にならないよう本文の行間を少し確保する．
\linespread{1.13}
\setlength{\parindent}{1zw}
\setlength{\parskip}{0pt}
\setlength{\abovedisplayskip}{6pt plus 2pt minus 2pt}
\setlength{\belowdisplayskip}{6pt plus 2pt minus 2pt}
\setlength{\abovedisplayshortskip}{3pt plus 2pt minus 1pt}
\setlength{\belowdisplayshortskip}{4pt plus 2pt minus 1pt}
\setlist[enumerate]{topsep=3pt,itemsep=1pt,parsep=0pt,leftmargin=2.4zw}
\titleformat{\section}{\normalfont\normalsize\bfseries}{\thesection}{0.8zw}{}
\titlespacing*{\section}{0pt}{2.0ex plus .4ex}{.9ex}
\titleformat{\subsection}{\normalfont\normalsize\bfseries}{\thesubsection}{0.7zw}{}
\titlespacing*{\subsection}{0pt}{1.5ex plus .3ex}{.6ex}
\raggedbottom

% ----- 柱とページ番号：ページ上部、罫線なし -----
\fancypagestyle{mathnote}{%
  \fancyhf{}%
  \fancyhead[L]{\footnotesize \NoteSeries\ （\NoteTitle）}%
  \fancyhead[R]{\footnotesize \thepage}%
  \renewcommand{\headrulewidth}{0pt}%
  \renewcommand{\footrulewidth}{0pt}%
}
\pagestyle{mathnote}

% ----- 節内通し番号の定義・定理群 -----
\newtheoremstyle{mathnotetheorem}
 {4.5pt}{4.5pt}{\normalfont}{}%
 {\bfseries}{}{0.65zw}%
 {\thmname{#1}\thmnumber{ #2}\thmnote{（#3）}}
\theoremstyle{mathnotetheorem}
\newtheorem{theorem}{定理}[section]
\newtheorem{definition}[theorem]{定義}
\newtheorem{lemma}[theorem]{補題}
\newtheorem{proposition}[theorem]{命題}
\newtheorem{corollary}[theorem]{系}
\newtheorem{example}[theorem]{例}
\newtheorem{remark}[theorem]{注意}
\renewcommand{\proofname}{証明}
\renewcommand{\qedsymbol}{$\square$}
\makeatletter
\renewenvironment{proof}[1][\proofname]{%
  \par\pushQED{\qed}\normalfont\topsep4\p@\@plus2\p@\relax
  \trivlist\item[\hskip\labelsep\bfseries #1]\ignorespaces
}{\popQED\endtrivlist\@endpefalse}
\makeatother

% ----- 使うノートで適宜拡張する記号 -----
\newcommand{\R}{\mathbb{R}}
\newcommand{\norm}[1]{\lVert #1\rVert}
\newcommand{\abs}[1]{\lvert #1\rvert}

\begin{document}
\section*{読了後にできること}
本稿を読むための前提は，自由群，自由加群，普遍写像性．
本稿の具体的な到達課題は，自由対象の普遍性を実際の写像の一意性として検証する．

\section{普遍性の定式化}

本稿では集合，群，可換群，可換環上の加群の
自由構成を扱う．これらを特徴づける共通の性質を
\emph{普遍性}という．

\begin{definition}
$S$ を集合，$\mathcal C$ をある種類の代数的対象の圏とする．
$F$ を $\mathcal C$ の対象，$\eta:S\to|F|$ をその台集合への写像とする．
$(F,\eta)$ が $S$ 上の\emph{自由対象}であるとは，
$\mathcal C$ の任意の対象 $A$ と任意の写像
$f:S\to|A|$ に対して，ただ一つの射
$\widetilde f:F\to A$ が存在し，
\[
 \widetilde f\circ\eta=f
\]
を満たすことをいう．
ここで $|A|$ は対象 $A$ の台集合とする．
\end{definition}

\begin{theorem}[普遍性による一意性]\label{thm:universal-unique}
$(F,\eta)$，$(F',\eta')$ がともに同じ集合 $S$ 上の
自由対象なら，ただ一つの同型
$\varphi:F\to F'$ で
$\varphi\circ\eta=\eta'$ となるものが存在する．
\end{theorem}
\begin{proof}
$F$ の普遍性を $A=F'$，$f=\eta'$ に適用すると，
ただ一つの射 $\varphi:F\to F'$ が存在して
$\varphi\eta=\eta'$ となる．
同様に $F'$ の普遍性から
$\psi:F'\to F$ で $\psi\eta'=\eta$ となるものが存在する．
このとき
\[
 (\psi\varphi)\eta=\psi\eta'=\eta
\]
だから，$F$ の普遍性によって
$\psi\varphi=\operatorname{id}_F$ である．
同様に $\varphi\psi=\operatorname{id}_{F'}$ となる．
従って $\varphi$ は同型であり，
その一意性も初めの普遍性から従う．
\end{proof}

\section{自由群の構成}

\begin{definition}
集合 $S$ に対し，
$S^\pm=S\times\{+1,-1\}$ を記号の集合とする．
有限列 $(s_1,\varepsilon_1)\cdots(s_n,\varepsilon_n)$ を
\emph{語}という．長さ零の語を空語と書く．
語全体は連結を積とするモノイド $W(S)$ をなす．
二つの連続した互いに逆の記号
$(s,+1)(s,-1)$ または $(s,-1)(s,+1)$ を
消去または挿入する操作で生成される
モノイド合同関係を $\sim$ とする．
\[
 F(S)=W(S)/{\sim}
\]
と置く．
\end{definition}

\begin{lemma}\label{lem:free-group}
$F(S)$ は群である．各 $s\in S$ の一文字語の同値類を
$\eta(s)$ とすると，
$\eta:S\to F(S)$ が定まる．
\end{lemma}
\begin{proof}
$\sim$ はモノイド合同なので，連結は商集合上に
結合的な積を定める．空語の同値類は単位元である．
語 $w=(s_1,\varepsilon_1)\cdots(s_n,\varepsilon_n)$ に対して
\[
 w^{-1}=(s_n,-\varepsilon_n)\cdots(s_1,-\varepsilon_1)
\]
と置く．$ww^{-1}$ および $w^{-1}w$ は
中央から隣接した逆の記号を順次消去することで
空語と同値になる．
従って任意の同値類が両側逆元を持つので
$F(S)$ は群である．
\end{proof}

\begin{theorem}[自由群の普遍性]\label{thm:free-group}
$S$ を集合とする．
$(F(S),\eta)$ は群の圏で $S$ 上の自由対象である．
\end{theorem}
\begin{proof}
任意の群 $G$ と写像 $f:S\to G$ を取る．
語 $w=(s_1,\varepsilon_1)\cdots(s_n,\varepsilon_n)$ に対し
\[
 \widehat f(w)=
 f(s_1)^{\varepsilon_1}\cdots f(s_n)^{\varepsilon_n}
\]
と定め，空語は $1_G$ に写す．
$\widehat f$ は語の連結を群の積に写すモノイド準同型である．
隣接した逆の記号の挿入・消去は，
像において $f(s)f(s)^{-1}=1_G$ を挿入・消去する操作なので，
$\widehat f$ は同値な語に同じ値を与える．
従って商集合上に群準同型
$\widetilde f:F(S)\to G$ を誘導する．
一文字語について
$\widetilde f(\eta(s))=f(s)$ なので
$\widetilde f\eta=f$ である．
一方，$F(S)$ の全ての元は $\eta(s)$ とその逆元の
有限積として表されるので，
$\eta(s)$ 上の値が定まれば群準同型の値は全て決まる．
従って拡張は一意である．
\end{proof}

\begin{proposition}[生成元の関係を課すこと]\label{prop:presentation}
集合 $S$ と自由群 $F(S)$，
関係式の集合 $R\subset F(S)$ に対し，
$N$ を $R$ が生成する正規部分群とし，
$G=F(S)/N$ と置く．
任意の群 $H$ と写像 $f:S\to H$ について，
$f$ が $G\to H$ の準同型へ拡張されることと，
全ての $r\in R$ に対して
$\widetilde f(r)=1_H$ が成立することは同値である．
ここで $\widetilde f:F(S)\to H$ は自由群の普遍性による写像である．
\end{proposition}
\begin{proof}
$G\to H$ への拡張があるなら，商写像
$\pi:F(S)\to G$ の核 $N$ は合成写像の核に含まれるため
全ての関係式は $1_H$ に写る．
逆に全ての $r\in R$ が $\widetilde f$ の核に属すれば，
核は正規部分群なので $N\subset\ker\widetilde f$ となる．
商群の普遍性から $\widetilde f$ は
唯一の群準同型 $G\to H$ を誘導する．
\end{proof}

\section{自由Abel群}

\begin{definition}
集合 $S$ と可換群 $A$ に対し，
\[
 A^{(S)}
 =\{u:S\to A:\{s\in S:u(s)\ne0\}\text{ は有限}\}
\]
と置く．点ごとの加法によって可換群をなす．
特に $\mathbb Z^{(S)}$ を $S$ 上の自由Abel群という．
$s\in S$ に対して
$\delta_s(t)=1$（$t=s$），$0$（$t\ne s$）
と定める．
\end{definition}

\begin{theorem}[自由Abel群の普遍性]\label{thm:free-abelian}
任意の可換群 $A$ と写像 $f:S\to A$ に対し，
ただ一つの群準同型
$\widetilde f:\mathbb Z^{(S)}\to A$ が存在して
$\widetilde f(\delta_s)=f(s)$（$s\in S$）となる．
\end{theorem}
\begin{proof}
$u\in\mathbb Z^{(S)}$ は有限台を持つので
\[
 \widetilde f(u)=\sum_{s\in S}u(s)f(s)
\]
と定めれば，有限和として意味を持つ．
点ごとの加法に関して
$\widetilde f(u+v)=\widetilde f(u)+\widetilde f(v)$
が成立するから群準同型である．
$\delta_s$ については
$\widetilde f(\delta_s)=f(s)$ となる．
任意の $u$ は
$u=\sum_{s\in S}u(s)\delta_s$ と一意に書けるので，
準同型の $\delta_s$ 上の値が決まれば
全ての $u$ の像が決まり，一意性が従う．
\end{proof}

\begin{remark}
自由群 $F(S)$ からその交換子部分群で割った商
$F(S)_{\mathrm{ab}}=F(S)/[F(S),F(S)]$ は，
自由Abel群 $\mathbb Z^{(S)}$ と自然に同型となる．
\end{remark}
\begin{proof}
任意の可換群 $A$ への群準同型は
交換子を零へ写すから，
群の自由性と商群の普遍性によって
$F(S)_{\mathrm{ab}}$ は可換群における $S$ 上の
自由対象となる．
定理\ref{thm:universal-unique}と
定理\ref{thm:free-abelian}によって一意な同型が存在する．
\end{proof}

\section{自由加群とその函手性}

以下，$R$ を単位元を持つ可換環とする．

\begin{definition}
集合 $S$ に対して
\[
 R^{(S)}=\{u:S\to R:\operatorname{supp}(u)\text{ は有限}\}
\]
に点ごとの加法・$R$-スカラー倍を定める．
この $R$-加群を $S$ 上の自由 $R$-加群という．
\end{definition}

\begin{theorem}[自由加群の普遍性]\label{thm:free-module}
任意の $R$-加群 $M$ と写像 $f:S\to M$ に対し，
ただ一つの $R$-線形写像
$\widetilde f:R^{(S)}\to M$ が存在して
$\widetilde f(\delta_s)=f(s)$ を満たす．
\end{theorem}
\begin{proof}
任意の $u\in R^{(S)}$ を
$u=\sum_{s\in S}u(s)\delta_s$ と一意に書ける．
そこで
$\widetilde f(u)=\sum_{s\in S}u(s)f(s)$ と定める．
有限和なので意味を持ち，線形性は
$R$-加群の公理から従う．
$\delta_s$ 上での値は $f(s)$ である．
逆に，$R$-線形写像は $\delta_s$ 上での値によって
全ての $u$ 上で決まるため一意である．
\end{proof}

\begin{proposition}[自由加群の函手性]\label{prop:free-functor}
集合写像 $a:S\to T$ に対し，
ただ一つの線形写像
$R^{(a)}:R^{(S)}\to R^{(T)}$ が存在して
$R^{(a)}(\delta_s)=\delta_{a(s)}$ となる．
さらに
\[
 R^{(\operatorname{id}_S)}
 =\operatorname{id}_{R^{(S)}},\qquad
 R^{(b\circ a)}=R^{(b)}\circ R^{(a)}
\]
が成り立つ．
\end{proposition}
\begin{proof}
最初の存在と一意性は
定理\ref{thm:free-module}に $M=R^{(T)}$，
$f(s)=\delta_{a(s)}$ を適用することで従う．
二つの函手性の等式について，
左右の線形写像は任意の基底ベクトル $\delta_s$ を
それぞれ同じベクトルに写す．
自由加群の普遍性から全ての元で一致する．
\end{proof}

\section{テンソル積の存在}

\begin{definition}
$R$-加群 $M,N,P$ に対し，
写像 $b:M\times N\to P$ が\emph{$R$-双線形}であるとは，
各変数を固定したとき残りの変数について
$R$-線形となることをいう．
\end{definition}

\begin{theorem}[テンソル積の存在と普遍性]\label{thm:tensor}
任意の $R$-加群 $M,N$ に対して，
$R$-加群 $M\otimes_RN$ と双線形写像
\[
 \tau:M\times N\to M\otimes_RN,\qquad
 (m,n)\longmapsto m\otimes n
\]
が存在し，任意の $R$-加群 $P$ と双線形写像
$b:M\times N\to P$ に対して，
ただ一つの $R$-線形写像
$\widetilde b:M\otimes_RN\to P$ が存在して
$b=\widetilde b\circ\tau$ を満たす．
\end{theorem}
\begin{proof}
集合 $M\times N$ 上の自由加群
$F=R^{(M\times N)}$ を取る．
$(m,n)$ に対応する基底ベクトルを $[m,n]$ と書く．
$F$ の部分加群 $L$ を，以下の形の元全体で生成される
部分加群とする：
\begin{align*}
 &[m+m',n]-[m,n]-[m',n],\\
 &[m,n+n']-[m,n]-[m,n'],\\
 &[rm,n]-r[m,n],\\
 &[m,rn]-r[m,n],
\end{align*}
ここで $m,m'\in M$，$n,n'\in N$，$r\in R$ とする．
\[
 M\otimes_RN=F/L
\]
と定め，$m\otimes n=[m,n]+L$ と置く．
$L$ の生成関係から
$\tau$ が各変数について $R$-線形であることが従う．

任意の双線形写像 $b:M\times N\to P$ を取る．
自由加群の普遍性により
$F$ 上の唯一の線形写像
$\widehat b:F\to P$ で
$\widehat b([m,n])=b(m,n)$ となるものが存在する．
双線形性から，$L$ の各生成元は
$\widehat b$ の核に属する．従って $L\subset\ker\widehat b$
であり，商加群の普遍性によって
ただ一つの線形写像
$\widetilde b:F/L\to P$ が誘導される．
これは
$\widetilde b(m\otimes n)=b(m,n)$ を満たす．
最後に $F/L$ は $m\otimes n$ の有限線形結合で
生成されるので，この値を固定した線形写像は一意である．
\end{proof}

\begin{proposition}[テンソル積の一意性]\label{prop:tensor-unique}
$R$-加群 $T$ と双線形写像
$\tau':M\times N\to T$ が
定理\ref{thm:tensor}と同じ普遍性を満たすなら，
ただ一つの $R$-加群同型
$\varphi:M\otimes_RN\to T$ が存在して
$\varphi(m\otimes n)=\tau'(m,n)$ を満たす．
\end{proposition}
\begin{proof}
$\tau'$ の普遍性と $\tau$ の普遍性をそれぞれ使うと，
$\varphi:M\otimes_RN\to T$，
$\psi:T\to M\otimes_RN$ で
$\varphi\tau=\tau'$，$\psi\tau'=\tau$ を満たす
線形写像が一意に存在する．
$\psi\varphi$ と恒等写像は $\tau$ との合成が一致するので，
普遍性により $\psi\varphi=\operatorname{id}$ である．
同様に $\varphi\psi=\operatorname{id}$ を得る．
従って $\varphi$ は同型であり，一意である．
\end{proof}

\section{テンソル積の基本性質}

\begin{proposition}[函手性]\label{prop:tensor-functor}
$R$-線形写像 $f:M\to M'$，
$g:N\to N'$ に対して，ただ一つの線形写像
\[
 f\otimes g:M\otimes_RN\to M'\otimes_RN'
\]
が存在して
$(f\otimes g)(m\otimes n)=f(m)\otimes g(n)$ となる．
また
\[
 \operatorname{id}_M\otimes\operatorname{id}_N
 =\operatorname{id}_{M\otimes_RN},\qquad
 (f'\otimes g')\circ(f\otimes g)
 =(f'\circ f)\otimes(g'\circ g).
\]
\end{proposition}
\begin{proof}
写像
$(m,n)\mapsto f(m)\otimes g(n)$ は双線形だから，
定理\ref{thm:tensor}によって
唯一の線形写像 $f\otimes g$ を誘導する．
二つの函手性等式は，それぞれ
任意の純テンソル $m\otimes n$ 上で左右が一致することを
確認し，純テンソルが全体を生成することから従う．
\end{proof}

\begin{theorem}[対称性と結合性]\label{thm:tensor-basic}
$R$-加群 $M,N,P$ に対して自然な同型
\[
 M\otimes_RN\cong N\otimes_RM,\qquad
 (M\otimes_RN)\otimes_RP\cong M\otimes_R(N\otimes_RP)
\]
が存在する．また
\[
 R\otimes_RM\cong M
\]
が成り立つ．
\end{theorem}
\begin{proof}
対称性について，
$(m,n)\mapsto n\otimes m$ は双線形なので，
$\sigma:M\otimes N\to N\otimes M$ が一意に存在して
$\sigma(m\otimes n)=n\otimes m$ となる．
逆向きも同じ方法で構成すると，
合成は任意の純テンソルで恒等写像に一致するので
互いに逆となる．

結合性について，三重線形写像
\[
 (m,n,p)\longmapsto m\otimes(n\otimes p)
\]
を考える．固定した $p$ に対して
$(m,n)\mapsto m\otimes(n\otimes p)$ は双線形であるため，
$M\otimes N$ 上の線形写像を誘導する．
これを $p$ に関して変化させると，
$(u,p)\mapsto u\mapsto m\otimes(n\otimes p)$ に対応する
双線形写像
$(M\otimes N)\times P\to M\otimes(N\otimes P)$ が得られる．
よって
\[
 \alpha:(M\otimes N)\otimes P\to M\otimes(N\otimes P)
\]
で
$\alpha((m\otimes n)\otimes p)=m\otimes(n\otimes p)$
となるものが存在する．
逆向きは
$(m,n,p)\mapsto(m\otimes n)\otimes p$ を用いて同様に構成する．
合成は全ての三重純テンソルで恒等写像に等しいから，
両者は同型である．

最後に
$\mu:R\times M\to M$，
$\mu(r,m)=rm$ は双線形なので
$\overline\mu:R\otimes M\to M$ を誘導する．
逆向きに $m\mapsto1\otimes m$ と置くと，
これは線形であり，
$\overline\mu(1\otimes m)=m$ かつ
\[
 1\otimes(rm)=r\otimes m
\]
がテンソル積の関係式から従う．
従って合成は両方向とも恒等写像である．
\end{proof}

\begin{theorem}[直和との可換性]\label{thm:tensor-sum}
$R$-加群 $M,N,P$ に対して
\[
 M\otimes_R(N\oplus P)
 \cong(M\otimes_RN)\oplus(M\otimes_RP)
\]
という自然な同型が存在する．
\end{theorem}
\begin{proof}
双線形写像
\[
 M\times(N\oplus P)\to
 (M\otimes N)\oplus(M\otimes P),\quad
 (m,(n,p))\mapsto(m\otimes n,m\otimes p)
\]
から線形写像 $\Phi$ を得る．
逆向きには，包含
$N\to N\oplus P$，$n\mapsto(n,0)$ と
$P\to N\oplus P$，$p\mapsto(0,p)$ から
\[
 \Psi:(M\otimes N)\oplus(M\otimes P)
 \to M\otimes(N\oplus P)
\]
を構成し，
$\Psi(m\otimes n,0)=m\otimes(n,0)$，
$\Psi(0,m\otimes p)=m\otimes(0,p)$ と定める．
$\Phi\Psi$ は直和の各成分上で恒等写像であり，
\[
 \Psi\Phi(m\otimes(n,p))
 =m\otimes(n,0)+m\otimes(0,p)
 =m\otimes(n,p)
\]
だから $\Psi\Phi$ も恒等写像である．
\end{proof}

\section{テンソル積とHomの随伴}

\begin{theorem}[テンソル--Homの随伴]\label{thm:tensor-hom}
$R$ を可換環，$M,N,P$ を $R$-加群とする．
このとき自然な $R$-加群同型
\[
 \operatorname{Hom}_R(M\otimes_RN,P)
 \cong\operatorname{Hom}_R(M,\operatorname{Hom}_R(N,P))
\]
が存在する．
\end{theorem}
\begin{proof}
$F:M\otimes N\to P$ を線形写像とする．
各 $m\in M$ に対して
$\widehat F(m):N\to P$ を
\[
 \widehat F(m)(n)=F(m\otimes n)
\]
と定める．テンソル積の双線形性と $F$ の線形性により
$\widehat F(m)$ は $R$-線形である．
さらに $m\mapsto\widehat F(m)$ も線形であるため，
\[
 \widehat F\in
 \operatorname{Hom}_R(M,\operatorname{Hom}_R(N,P))
\]
となる．
逆に $G:M\to\operatorname{Hom}_R(N,P)$ が線形なら，
\[
 b_G:M\times N\to P,\qquad b_G(m,n)=G(m)(n)
\]
は双線形である．テンソル積の普遍性によって
ただ一つの線形写像
$\widetilde b_G:M\otimes N\to P$ が存在し，
$\widetilde b_G(m\otimes n)=G(m)(n)$ となる．
二つの対応を合成すると，
全ての純テンソルおよび全ての $m,n$ に対して元に戻るので，
互いに逆である．
また写像の和と $R$-スカラー倍に関して両対応は線形である．
従って自然な $R$-加群同型を得る．
\end{proof}

\begin{example}[有限自由加群のテンソル積]
$R^m\otimes_RR^n\cong R^{mn}$ である．
\end{example}
\begin{proof}
$R^m$ の標準基底を $e_i$，
$R^n$ の標準基底を $f_j$ とする．
任意の純テンソルは
\[
 \left(\sum_i a_ie_i\right)\otimes
 \left(\sum_j b_jf_j\right)
 =\sum_{i,j}a_ib_j(e_i\otimes f_j)
\]
と表されるので $e_i\otimes f_j$ が全体を生成する．
一方，双線形写像
$(x,y)\mapsto(x_iy_j)_{i,j}$ は
線形写像 $R^m\otimes R^n\to R^{mn}$ を誘導する．
これは $e_i\otimes f_j$ を標準基底へ写すから，
それらが独立であることも従う．
従って同型である．
\end{proof}

\section{自由--忘却の随伴}

\begin{theorem}[自由構成の随伴性]\label{thm:free-forgetful}
集合 $S$ と $R$-加群 $M$ に対して，
自然な全単射
\[
 \operatorname{Hom}_R(R^{(S)},M)
 \cong\operatorname{Map}(S,|M|)
\]
が存在する．この対応は，集合写像 $S'\to S$ と
線形写像 $M\to M'$ のいずれに関しても自然である．
\end{theorem}
\begin{proof}
線形写像 $F:R^{(S)}\to M$ を
$F\circ\eta:S\to|M|$ に対応させる．
定理\ref{thm:free-module}により，
任意の $f:S\to|M|$ に対して
ただ一つの線形写像 $\widetilde f:R^{(S)}\to M$ が存在する．
二つの対応は互いに逆であるから全単射となる．

自然性を示す．$a:S'\to S$，$G:M\to M'$ を取る．
任意の $F:R^{(S)}\to M$ に対して
\[
 (G\circ F\circ R^{(a)})\eta_{S'}
 =G\circ F\circ\eta_S\circ a.
\]
これは左辺の線形写像と右辺の集合写像が，
対応を通じて一致することを意味する．
従って自然性が成り立つ．
\end{proof}

\begin{remark}
自由群，自由Abel群，自由加群，テンソル積は，
いずれも普遍性によって定義することができる．
普遍性を用いる利点は，具体的構成に依存しない一意性と，
写像・同型・函手性が同じ形式の議論から導かれることである．
\end{remark}


\section{到達課題の完全解答}
\noindent\textbf{自由Abel群の普遍性．}
\begin{proposition}
集合 $S=\{a,b\}$ 上の自由Abel群は
$\mathbb Z^2$ であり，任意のAbel群 $A$ と
写像 $f:S\to A$ に対して一意な準同型
$\widetilde f:\mathbb Z^2\to A$ が存在して
$\widetilde f(1,0)=f(a)$，
$\widetilde f(0,1)=f(b)$ を満たす．
\end{proposition}
\begin{proof}
$\widetilde f(m,n)=mf(a)+nf(b)$ と定めれば
加法性を満たす．
$\mathbb Z^2$ の元は
$(m,n)=m(1,0)+n(0,1)$ と一意に書けるので，
指定された二元の像を持つ準同型はこの式以外にない．
\end{proof}
\begin{corollary}
$\operatorname{Hom}_{\mathbb Z}(\mathbb Z^2,A)
\cong A\times A$ である．
\end{corollary}
\begin{proof}
準同型を $(1,0),(0,1)$ の像へ送れば，
上の普遍性によって全単射となる．
\end{proof}
\noindent\textbf{到達確認．}
$S$ が3元集合なら
$S$ 上の自由Abel群から $\mathbb Z/2\mathbb Z$ への
準同型は何個あるか．
\par\noindent\textbf{解答．}
各生成元の像は独立に2通り選べるため
$2^3=8$ 個である．
\end{document}